%Master File:lectures.tex
\lesson{Introduction to Lab 1}
\vspace{-2cm}
\begin{center}
  \includegraphics[height=10cm]{Cell_Cycle.png}
\end{center}

\keywords{Chemical Equations, Stochasticity, Gillespie Algorithm}
%%%%%%%%%%%%%%%%%%%%%%% Next Slide %%%%%%%%%%%%%%%%%%%%%%%
\renewcommand{\Outline}{%
\begin{slide}
\section[1]{Outline}

\begin{minipage}{12cm}
  \begin{enumerate}\squeeze
    \outlineitem{The Tutorials}{tutorial}
    \outlineitem{Lab 1}{labone}
   \end{enumerate}
\end{minipage}\hfill
\begin{minipage}{10cm}
  \vfill
  \includegraphics[width=10cm]{Cell_Cycle.png}
  \vfill
\end{minipage}
\end{slide}
\addtocounter{outlineitem}{1}
}

\setcounter{outlineitem}{1}
\Outline
\toptarget{firstoutline}

%%%%%%%%%%%%%%%%%%%%%%% Next Slide %%%%%%%%%%%%%%%%%%%%%%%

\begin{slide}
\section{The Labs/Tutorials}

\begin{PauseHighLight}
  \begin{itemize}
  \item There are four tutorials\pause
  \item The first tutorial suggests some videos for you to look
    at\pause---do this in your own time\pauseb
  \item The first and later tutorials describes works you need to
    do\pause
  \item The coursework is to be done individually\pause
  \item Each tutorial is worth 7.5\%\pause
  \end{itemize}
\end{PauseHighLight}

\end{slide}
%%%%%%%%%%%%%%%%%%%%%%% Next Slide %%%%%%%%%%%%%%%%%%%%%%%

\begin{slide}
  \section{Hand-In}
  
\begin{PauseHighLight}
  \begin{itemize}
  \item Use the hand-in system to submit your work as a PDF (not as a
    jupyter notebook)\pause
  \item Hand-in dates are
    \begin{enumerate}
    \item Due 4pm on Thursday 18th February 2022\pause
    \item Due 4pm on Thursday 24th March 2022
    \item Due 4pm on Thursday 5th May 2022
    \item Due 4pm on Thursday 19th May 2022\pause
    \end{enumerate}
  \end{itemize}
\end{PauseHighLight}

\end{slide}

%%%%%%%%%%%%%%%%%%%%%%% Next Slide %%%%%%%%%%%%%%%%%%%%%%%

\begin{slide}
\section{Help}

\begin{PauseHighLight}
  \begin{itemize}
  \item We will run this session every week\pause
  \item We will discuss the labs and broad issues that might help
    everyone\pause
  \item We will then split into channels (hence you groups)---I
    intended to split you into groups of 5, but because of the limitations
    of Teams I have put two groups together\pause
  \item If you are not yet in a group the just join a channel\pause
  \item A team of four of us will go around the 10 channels\pause
  \end{itemize}
\end{PauseHighLight}

\end{slide}

%%%%%%%%%%%%%%%%%%%%%%% Next Slide %%%%%%%%%%%%%%%%%%%%%%%
\Outline
%%%%%%%%%%%%%%%%%%%%%%% Next Slide %%%%%%%%%%%%%%%%%%%%%%%

\begin{slide}
\section{Data Exploration}

\begin{PauseHighLight}
  \begin{itemize}
  \item The first lab looks at biological data exploration\pause
  \item The data consists of RNA measurements and raw protein
    measurements\pause
  \item RNA measurements acts as a proxy for protein concentration (it
    is cheap to measure)\pause
  \item The protein measurements are real protein concentration but
    much more expensive to do\pause
  \end{itemize}
\end{PauseHighLight}

\end{slide}

%%%%%%%%%%%%%%%%%%%%%%% Next Slide %%%%%%%%%%%%%%%%%%%%%%%

\begin{slide}
\section{The Data}

\begin{PauseHighLight}
  \begin{itemize}
  \item The data is from humans\pause
  \item It is measure at three points in the cell-cycle\pause
  \item That is, cells regularly grow and divide\pause
  \item These cells were cleverly synchronised so we could measure the
    protein concentration at different points in the cell cycle\pause
  \end{itemize}
\end{PauseHighLight}

\end{slide}

%%%%%%%%%%%%%%%%%%%%%%% Next Slide %%%%%%%%%%%%%%%%%%%%%%%

\begin{slide}
\section[-2]{The Cell Cycle}

\begin{center}
  \includegraphics[width=0.6\linewidth]{Cell_Cycle.png}
\end{center}
\end{slide}

%%%%%%%%%%%%%%%%%%%%%%% Next Slide %%%%%%%%%%%%%%%%%%%%%%%

\begin{slide}
  \section{Python and Jupyter Notebook}

\begin{PauseHighLight}
  \begin{itemize}
  \item We provide python code to do the basic data exploration\pause
  \item You are welcome to use this or to ignore it\pause
  \item To get full marks we are expecting you to have gone beyond what
    we give you\pause
  \item We provide a jupyter notebook.  You don't need to use this, but
    it is a reasonably nice way to visualise what you are doing\pause
  \item Please don't submit the raw notebook---you can turn this into PDF\pause
  \end{itemize}
\end{PauseHighLight}

\end{slide}

%%%%%%%%%%%%%%%%%%%%%%% Next Slide %%%%%%%%%%%%%%%%%%%%%%%

\begin{slide}
\section{Pandas}

\begin{PauseHighLight}
  \begin{itemize}
  \item Pandas is a python package that dominates modern data science\pause
  \item Other package are available\pause---but they are nearly all
    based on pandas\pauseb
  \item Pandas has two data structures: \texttt{Series} and
    \texttt{DataFrame}\pause
  \item DataFrame is a glorified table with labels for the columns and
    rows\pause
  \item It has a lot of functionality built around it for rapid data exploration\pause
  \end{itemize}
\end{PauseHighLight}

\end{slide}

%%%%%%%%%%%%%%%%%%%%%%% Next Slide %%%%%%%%%%%%%%%%%%%%%%%

\begin{slide}
\section{Matplotlib}

\begin{PauseHighLight}
  \begin{itemize}
  \item Matplotlib is a visualisation library\pause
  \item It was designed to look like MATLAB plot functions\pause{} and
    is consequently very ugly\pauseb
  \item It has a lot of functionality, but isn't particularly natural
    to use\pause
  \item Seaborn is build on top of matplotlib and produces slightly
    prettier graphs (but is still not easy to use)\pause
  \item Sometime you need to write \texttt{> \%matplotlib inline} to
    get the graphs to show\pause
  \end{itemize}
\end{PauseHighLight}

\end{slide}
